\documentclass[a4paper,12pt]{article}
\usepackage[T2A]{fontenc}
\usepackage[utf8]{inputenc}
\usepackage[russian]{babel}
\usepackage{amsmath,amssymb,mathtools}
\usepackage{PTSans}
\renewcommand{\familydefault}{\sfdefault}
\usepackage{geometry}
\geometry{left=23mm,right=23mm,top=24mm,bottom=25mm}
\usepackage{microtype}
\usepackage{tikz}
\usepackage{booktabs,tabularx}
\usepackage{xcolor}
\usepackage{hyperref}
\usepackage{fancyhdr}
\usepackage{setspace}
\definecolor{accent}{HTML}{13777A}
\definecolor{darkaccent}{HTML}{135153}
\definecolor{textgray}{HTML}{373B3B}
\hypersetup{colorlinks=true,linkcolor=darkaccent,pdfauthor={Лёвин Артём Александрович},pdftitle={Ревью домашней работы — Тимофей Васильченко}}
\setlength{\parindent}{0pt}
\setlength{\parskip}{3pt}
\setstretch{1.10}
\setlength{\headheight}{15pt}
\pagestyle{fancy}
\fancyhf{}
\fancyhead[L]{\small\textcolor{darkaccent}{Ревью домашней работы}}
\fancyhead[R]{\small\textcolor{textgray}{31 октября 2026 г.}}
\fancyfoot[C]{\scriptsize\textcolor{textgray}{Лёвин Артём Александрович эксклюзивно для Тимофея Васильченко}}
\fancyfoot[R]{\small\thepage}
\renewcommand{\headrulewidth}{0.3pt}
\renewcommand{\footrulewidth}{0pt}
\newcommand{\headtxt}[1]{\par\smallskip\textbf{\textcolor{darkaccent}{#1}}\par\vspace{-1mm}}
\begin{document}
\thispagestyle{empty}
\begin{center}
\vspace*{39mm}
{\large\textcolor{darkaccent}{МАТЕМАТИКА}}\par
\vspace{18mm}
{\Huge\bfseries Ревью домашней работы}\par
\vspace{12mm}
{\large Подготовка к ЕГЭ по профильной математике}\par
\vfill
{\large Лёвин Артём Александрович}\par
\vspace{6mm}
{\large 31 октября 2026 г.}\par
\vspace{25mm}
\begin{tikzpicture}
\draw[accent!60,line width=0.9pt] (-5,0) .. controls (-2,0.7) and (2,-0.7) .. (5,0);
\end{tikzpicture}
\end{center}
\newpage
\section*{Сводная таблица}
\textcolor{textgray}{В таблице приведены только задания, в которых требуется исправление или отсутствует решение.}\par
\vspace{5mm}
\footnotesize
\setlength{\tabcolsep}{3pt}
\renewcommand{\arraystretch}{1.35}
\begin{tabular}{@{}p{8mm}p{24mm}p{60mm}p{24mm}p{31mm}@{}}
\toprule
\textbf{№} & \textbf{Тема} & \textbf{Суть замечания} & \textbf{Ваш ответ} & \textbf{Верный ответ}\\
\midrule
\hyperref[task:12]{12} & Логарифмы & Решение отсутствует. & --- & $\dfrac18$\\
\hyperref[task:13]{13} & Кредит & Лишнее деление на 2; округление ответа. & $2{,}06$ млн руб. & $2\,062\,500$ руб.\\
\bottomrule
\end{tabular}\par
\normalsize
\vfill
\textcolor{textgray}{Разборы доступны по нажатию на номера заданий в таблице.}
\newpage
\thispagestyle{empty}
\null\vfill
\begin{center}
{\Huge\bfseries\textcolor{darkaccent}{Разбор ошибок}}
\end{center}
\vfill\null
\newpage
\phantomsection\label{task:12}
\section*{Задание № 12. Логарифмическая функция}
\headtxt{Условие}
На графике изображена функция $f(x)=k\log_a x$. Найдите значение $x$, при котором $f(x)=-9$. На графике отмечена точка $(2;3)$.

\headtxt{Ваше решение и ответ}
В представленной работе решение и ответ к заданию № 12 отсутствуют.

\headtxt{\textcolor{red}{Ошибка:} отсутствует решение}
Последнего правильного вычислительного шага нет, поэтому указать первый неверный переход невозможно. Начать следует с чтения координат отмеченной точки на графике.

\headtxt{Почему это важно}
Число $a$ и коэффициент $k$ неизвестны по отдельности. Однако значение функции при $x=2$ известно. Этого достаточно, чтобы найти значение аргумента, при котором функция равна $-9$.

\headtxt{\textcolor{blue!60!black}{Ключевая идея:} свойство логарифма степени}
При $x>0$, $a>0$, $a\ne1$ верно
\[
\log_a(x^n)=n\log_a x.
\]
Поэтому для данной функции $f(2^n)=n f(2)$.

\headtxt{Полное правильное решение}
\textbf{Шаг 1.} По отмеченной точке графика считываем $f(2)=3$, то есть
\[
k\log_a2=3.
\]
\textbf{Шаг 2.} Требуется получить значение функции, противоположное утроенному числу~3:
\[
-9=-3\cdot3=-3k\log_a2.
\]
\textbf{Шаг 3.} Переносим множитель $-3$ в показатель логарифма:
\[
-3k\log_a2=k\log_a(2^{-3})=f(2^{-3}).
\]
Следовательно, $x=2^{-3}=\dfrac18$.

\headtxt{Проверка}
\[
f\!\left(\frac18\right)=k\log_a(2^{-3})=-3k\log_a2=-3\cdot3=-9.
\]
На данном графике функция возрастает, поэтому найденное значение $x$ единственно.

\headtxt{\textcolor{teal!60!black}{Итог:} $\displaystyle x=\frac18$}
\headtxt{Задача на закрепление}
Для функции $g(x)=k\log_a x$ известно, что $g(3)=2$. Найдите $x$, при котором $g(x)=-6$.

\newpage
\phantomsection\label{task:13}
\section*{Задание № 13. Кредит и равные уменьшения долга}
\headtxt{Условие}
В июле 2030 года взят кредит $1{,}5$ млн рублей на четыре года. В январе каждого года долг увеличивается на $15\%$ от остатка на конец предыдущего года. С февраля по июнь вносятся платежи; к июлю каждого следующего года долг уменьшается на одну и ту же величину. К июлю 2034 года долг полностью погашен. Найдите общую сумму выплат \textbf{в рублях}.

\headtxt{Ваше решение и ответ}
Вы верно обозначили $S=1{,}5$ млн, $p=0{,}15$, $N=4$ и записали формулу
\[
S+Sp\,\frac{N+1}{2}.
\]
Затем нашли $Sp=0{,}225$ млн и записали
\[
1{,}5+\frac{0{,}225\cdot2{,}5}{2}=0{,}56+1{,}5=2{,}06.
\]
Итоговый ответ в тетради: $2{,}06$.

\headtxt{\textcolor{red}{Ошибка:} лишнее деление на 2 и округление}
\textbf{Последний правильный шаг:} $S+Sp\dfrac{N+1}{2}$.

\textbf{Первый неверный шаг:} после подстановки $N=4$ выражение $\dfrac{N+1}{2}$ уже равно $\dfrac52=2{,}5$. Повторное деление на~2 в записи $0{,}225\cdot2{,}5/2$ лишнее.

Кроме того, $0{,}225\cdot2{,}5/2=0{,}28125$, а не $0{,}56$. Значение $0{,}56$ получается при вычислении $0{,}225\cdot2{,}5$ с округлением, но это округление теряет тысячи рублей.

\headtxt{\textcolor{blue!60!black}{Ключевая идея:} проценты на остаток долга}
Основной долг убывает равными долями, а проценты каждый год начисляются на ещё не погашенную сумму. Поэтому в четыре года проценты начисляются соответственно на $1$, $\dfrac34$, $\dfrac12$ и $\dfrac14$ первоначального долга.

\headtxt{Полное правильное решение}
\textbf{Шаг 1.} Ежегодное уменьшение основного долга:
\[
\frac{1\,500\,000}{4}=375\,000\ \text{руб.}
\]
\textbf{Шаг 2.} Остатки долга перед начислением процентов: $1\,500\,000$; $1\,125\,000$; $750\,000$; $375\,000$ рублей.

\textbf{Шаг 3.} Сумма всех начисленных процентов:
\begin{align*}
P&=0{,}15(1\,500\,000+1\,125\,000\\
 &\qquad\quad+750\,000+375\,000)\\
 &=0{,}15\cdot3\,750\,000=562\,500\ \text{руб.}
\end{align*}
\textbf{Шаг 4.} Общая сумма выплат состоит из первоначального долга и процентов:
\[
1\,500\,000+562\,500=2\,062\,500\ \text{руб.}
\]

\headtxt{Исправление по Вашей формуле}
\[
S+Sp\frac{N+1}{2}=1{,}5+0{,}225\cdot\frac52
=2{,}0625\ \text{млн руб.}
\]
Переводим миллионы в рубли: $2{,}0625\cdot1\,000\,000=2\,062\,500$.

\headtxt{Проверка}
Проценты за годы: $225\,000+168\,750+112\,500+56\,250=562\,500$ рублей. Результат совпал с расчётом по формуле.

\headtxt{\textcolor{teal!60!black}{Итог:} $2\,062\,500$ рублей}
\headtxt{Задача на закрепление}
Кредит $1{,}2$ млн рублей взят на три года. В январе каждого года долг увеличивается на $10\%$ от остатка на конец предыдущего года. К июлю каждого следующего года основной долг уменьшается на одну и ту же сумму; к концу третьего года кредит погашен. Найдите общую сумму выплат в рублях.
\end{document}
